% !TeX root = ../tesis.tex
%%%%%%%%%%%%%%%%%%%%%%%%%%%%%%%%%%%%%%%%%%%%%%%%%%%%%%%%%%%%%%%%%%%
% CAPÍTULO 6 — ESCENARIO HIPOTÉTICO: ANILLO PERIFÉRICO DE ALTA
%              CAPACIDAD
%
% Issue #25 [C-8]: Demostración de replicabilidad de la metodología
% mediante la misma cadena de herramientas aplicada a un GTFS
% hipotético que incorpora corredores anillares.
%
% PREREQUISITO: Cap. 5 completado y valores base validados.
%
% Estructura objetivo:
%   §6.1  Justificación del Escenario y Supuestos de Diseño
%         §6.1.1  Criterios de diseño del anillo hipotético
%         §6.1.2  Modalidad: Metrobús o Metro
%         §6.1.3  Cuatro cuadrantes del Anillo Periférico
%   §6.2  Construcción del GTFS Hipotético
%         §6.2.1  Metodología de generación del feed
%         §6.2.2  Estaciones, rutas, frecuencias y velocidades
%         §6.2.3  Nodos de transferencia con la red existente
%   §6.3  Procesamiento con la Cadena de Herramientas
%         §6.3.1  Ingestión en Apimetro
%         §6.3.2  Grafo ampliado en VFTModel
%         §6.3.3  Visualización en Transport-gis
%   §6.4  Indicadores de la Red Hipotética
%         §6.4.1  Cobertura ampliada (C)
%         §6.4.2  Fuerza Capilar con el anillo (FC_h)
%         §6.4.3  Factor de Desviación (DI)
%         §6.4.4  Tiempo de Viaje Promedio (T)
%         §6.4.5  Centralidad de Intermediación B(v)
%   §6.5  Comparación de Escenarios: Red Actual vs. Red con Anillo
%         §6.5.1  Tabla comparativa de indicadores
%         §6.5.2  Mapas de diferencias (delta)
%         §6.5.3  Zonas con mayor mejora en accesibilidad
%         §6.5.4  Interpretación urbanística
%%%%%%%%%%%%%%%%%%%%%%%%%%%%%%%%%%%%%%%%%%%%%%%%%%%%%%%%%%%%%%%%%%%
\chapter{Escenario Hipotético: Anillo Periférico de Alta Capacidad}
\label{ch:analisis-red-anillo}

% Secciones del capítulo (se agregan conforme se construya el GTFS
% hipotético y se ejecute la cadena de herramientas)
% \input{6-Analisis-Red-Anillo/6-1-Justificacion-Escenario}
% \input{6-Analisis-Red-Anillo/6-2-GTFS-Hipotetico}
% \input{6-Analisis-Red-Anillo/6-3-Procesamiento-Cadena}
% \input{6-Analisis-Red-Anillo/6-4-Indicadores-Hipoteticos}
% \input{6-Analisis-Red-Anillo/6-5-Comparacion-Escenarios}
